% !TEX root = ../main.tex
\hspace{24pt}

This chapter details each technical component of the proposed pipeline.
Section~\ref{sec:prerequisite} first introduces the prerequisite diffusion model architectures,
including the training principles of the DDPM denoising probabilistic model, the implicit deterministic sampling mechanism of DDIM,
the LDM latent diffusion architecture (with the VAE and the U-Net denoiser),
and the training and inference mechanisms of Marigold built upon this architecture,
serving as background knowledge for understanding the subsequent sparse-point-guided depth estimation of Marigold-DC.
Section~\ref{sec:depth} describes the depth estimation pipeline,
covering the acquisition of a rough depth map with Depth Anything V2\cite{depthanythingv2}, interactive sparse control point annotation,
and the \textit{ChinesePaintingLoss} loss function designed specifically for Traditional Chinese Painting.
Section~\ref{sec:pointcloud} introduces how the refined depth map is unprojected into a three-dimensional point cloud and rendered into a point-cloud video.
Section~\ref{sec:video} explains how the point-cloud video guides the PCDController of Uni3C to generate the final 2.5D video.
The system diagram of this study is shown in Figure~\ref{fig:pipeline}.

\section{Preliminaries: Diffusion Model Architectures}
\label{sec:prerequisite}

\subsection{Denoising Diffusion Probabilistic Models (DDPM)}

Denoising Diffusion Probabilistic Models (DDPM) \cite{ddpm}
are built around two mutually corresponding Markov chain processes.

\textbf{The forward process} is a fixed noising chain:
given the original data $\mathbf{x}_0$,
Gaussian noise is added sequentially at time steps $t = 1, \ldots, T$,
so that the data gradually approach pure noise:

\begin{equation}
  q(\mathbf{x}_t \mid \mathbf{x}_0) =
    \mathcal{N}\!\bigl(\mathbf{x}_t;\;\sqrt{\bar{\alpha}_t}\,\mathbf{x}_0,\;(1-\bar{\alpha}_t)\mathbf{I}\bigr)
\end{equation}

where $\bar{\alpha}_t = \prod_{s=1}^{t}(1 - \beta_s)$,
$\beta_s$ is a predefined noise schedule,
and $\mathbf{I}$ is the identity matrix, indicating that the noise is independent across dimensions.
When $T$ is sufficiently large, $\mathbf{x}_T \approx \mathcal{N}(\mathbf{0}, \mathbf{I})$.

\textbf{The reverse process} trains a neural network
$\boldsymbol{\epsilon}_\theta(\mathbf{x}_t, t)$ with trainable parameters $\theta$
to predict the noise added at each step,
using the mean squared error as the training objective:

\begin{equation}
  \mathcal{L}_{\mathrm{DDPM}} =
    \mathbb{E}_{\mathbf{x}_0,\,\boldsymbol{\epsilon},\,t}
    \bigl\|\boldsymbol{\epsilon} - \boldsymbol{\epsilon}_\theta(\mathbf{x}_t, t)\bigr\|^2
\end{equation}

where $\boldsymbol{\epsilon} \sim \mathcal{N}(\mathbf{0}, \mathbf{I})$
is the ground-truth noise sampled in the forward process, and $\boldsymbol{\epsilon}_\theta$ is the network prediction;
the difference between the two is the denoising error.

After training, the model starts from $\mathbf{x}_T \sim \mathcal{N}(\mathbf{0}, \mathbf{I})$;
at each step the neural network first estimates the denoised $\hat{\mathbf{x}}_0$,
and the state is then updated to $\mathbf{x}_{t-1}$ by stochastic sampling,
eventually recovering $\mathbf{x}_0$; the noising and denoising processes are illustrated in Figure~\ref{fig:ddpmrw}.

\subsection{Denoising Diffusion Implicit Models (DDIM)}

As shown in Figure~\ref{fig:ddpmrw}, the reverse sampling of DDPM denoises step by step along the Markov chain
$\mathbf{x}_T \to \cdots \to \mathbf{x}_t \to \mathbf{x}_{t-1} \to \cdots \to \mathbf{x}_0$,
where each step depends on the previous one and contains a stochastic noise term;
complete sampling requires $T$ (typically 1000) steps, making inference slow.
Denoising Diffusion Implicit Models (DDIM) \cite{ddim}
redefine the reverse sampling process as a non-Markovian one,
greatly accelerating inference while reusing the trained DDPM weights.

The key observation of DDIM stems from an analysis of the DDPM training procedure.
In the forward process, the joint distribution of the Markov chain is
$q(\mathbf{x}_{1:T} \mid \mathbf{x}_0) = \prod_{t=1}^T q(\mathbf{x}_t \mid \mathbf{x}_{t-1})$,
where each step depends on the previous one.
However, the DDPM training objective (the mean squared error loss) only requires sampling $\mathbf{x}_t$
directly from $\mathbf{x}_0$ at each time step $t$,
using the marginal distribution
$q(\mathbf{x}_t \mid \mathbf{x}_0) = \mathcal{N}(\sqrt{\bar{\alpha}_t}\,\mathbf{x}_0,\,(1-\bar{\alpha}_t)\mathbf{I})$.
The marginal distribution refers to the direct relation between $\mathbf{x}_t$ and $\mathbf{x}_0$
obtained by integrating out the intermediate variables
$\mathbf{x}_1, \ldots, \mathbf{x}_{t-1}$ of the joint distribution;
in DDPM this marginal has a closed-form solution, so $\mathbf{x}_t$ at any time step can be sampled from $\mathbf{x}_0$ in a single step.
DDIM observes that as long as a newly defined non-Markovian forward process shares the same marginal distribution
$q(\mathbf{x}_t \mid \mathbf{x}_0)$ as the original Markov chain,
the DDPM training loss is unaffected, and the trained network weights can therefore be reused directly.
Based on this, DDIM defines a fully deterministic update rule:
using the noise $\boldsymbol{\epsilon}_\theta(\mathbf{x}_t, t)$ predicted by the neural network,
it first recovers an estimate of the original data:

\begin{equation}
  \hat{\mathbf{x}}_0 =
    \frac{\mathbf{x}_t - \sqrt{1-\bar{\alpha}_t}\,\boldsymbol{\epsilon}_\theta(\mathbf{x}_t, t)}%
         {\sqrt{\bar{\alpha}_t}}
\end{equation}

and then updates to the previous time step with a deterministic formula:

\begin{equation}
  \mathbf{x}_{t-1} =
    \sqrt{\bar{\alpha}_{t-1}}\,\hat{\mathbf{x}}_0
    + \sqrt{1 - \bar{\alpha}_{t-1}}\,\boldsymbol{\epsilon}_\theta(\mathbf{x}_t, t)
\end{equation}

This update rule contains no stochastic noise term;
given the same initial noise $\mathbf{x}_T$, it necessarily produces the same $\mathbf{x}_0$.
Since each step does not depend on the randomness of the previous one,
DDIM can subsample the sequence of time steps,
skipping most of the intermediate steps
and completing in only 20--50 steps the sampling that originally required 1000 steps,
accelerating inference by roughly 10 to 50 times with almost no degradation in generation quality.
Moreover, deterministic sampling makes DDIM inversion possible:
forward noising of a known image can precisely recover its corresponding latent noise $\mathbf{x}_T$,
providing a key technical foundation for image editing and image-conditioned generation tasks.

\subsection{Latent Diffusion Models (LDM)}

Since both DDPM and DDIM operate in pixel space,
their memory and computation demands for high-resolution images are extremely high.
Latent Diffusion Models (LDM) \cite{ldm}
move the diffusion process into a low-dimensional latent space to reduce the computational cost;
the overall architecture, shown in Figure~\ref{fig:ldmrw}, consists of two main components.

\textbf{The Variational Autoencoder (VAE)}
is responsible for the conversion between pixel space and latent space;
its architecture originates from the VQGAN proposed by Esser et al. \cite{vqgan},
and LDM adopts the same network structure but replaces vector quantization (VQ) with KL-divergence regularization as the latent-space constraint.
The encoder $\mathcal{E}$ compresses an image $\mathbf{x} \in \mathbb{R}^{H \times W \times 3}$
into a latent vector $\mathbf{z} = \mathcal{E}(\mathbf{x}) \in \mathbb{R}^{h \times w \times c}$;
taking the Stable Diffusion v2 VAE adopted in this study as an example,
its spatial downsampling factor is $f=8$,
i.e., both height and width are reduced by a factor of 8 ($h = H/8$, $w = W/8$),
the total number of spatial pixels is compressed to $1/f^2 = 1/64$ of the original, and the number of channels is $c = 4$;
the decoder $\mathcal{D}$ maps the latent vector back to pixel space: $\hat{\mathbf{x}} = \mathcal{D}(\mathbf{z})$.
Natural images and depth maps exhibit substantial spatial redundancy in pixel space---
neighboring pixels are highly correlated semantically, and their true information content is far below the raw pixel count.
By jointly minimizing the reconstruction loss and
the KL-divergence regularization imposed on the latent distribution,
the VAE learns a low-dimensional latent manifold,
so that $\mathcal{E}$ can reduce each spatial dimension by a factor of $f=8$ (compressing the total pixel count by $f^2=64$)
while preserving the semantic and structural information sufficient to reconstruct the original content,
and $\mathcal{D}$ recovers the pixel details from this manifold.

\textbf{The U-Net denoiser} is the core computational module of LDM;
it is based on the U-Net proposed by Ronneberger et al. \cite{unet}
and is built as a symmetric encoder--decoder architecture with skip connections (the central green region in Figure~\ref{fig:ldmrw}).
The conditional encoder $\tau_\theta$ on the right accepts various forms of conditioning input
(semantic maps, text, feature representations, images),
and injects the encoded vectors into the layers of the U-Net
through the cross-attention mechanism (the Q/K/V modules in the figure),
so that the denoising prediction is guided by the conditioning content.

\subsection{Marigold}
\label{subsec:marigold}

The overall architecture of Marigold has been introduced in Chapter 2;
this section elaborates on its concrete training and inference mechanisms,
serving as the theoretical basis for the subsequent Marigold-DC guidance mechanism and the design of \textit{ChinesePaintingLoss}.
The overall training procedure is illustrated in Figure~\ref{fig:marigoldrw}.

\paragraph{Training}
Marigold takes an RGB image $\mathbf{x}$ and its corresponding depth map $\mathbf{d}$ as a training pair.
The VAE encoder (marked with a lock in the figure, indicating that its weights are frozen during training)
compresses both into the latent space:
$\mathbf{z}^{(\mathbf{x})} = \mathcal{E}(\mathbf{x})$ is the image latent code,
and $\mathbf{z}^{(\mathbf{d})} = \mathcal{E}(\mathbf{d})$ is the depth latent code.
A time step is sampled from the uniform distribution $t \sim \mathcal{U}(T)$, and noise is sampled from $\boldsymbol{\epsilon} \sim \mathcal{N}(\mathbf{0},\mathbf{I})$;
the depth latent code is then noised to obtain $\mathbf{z}_t^{(\mathbf{d})}$.
The image latent code $\mathbf{z}^{(\mathbf{x})}$ and the noisy depth $\mathbf{z}_t^{(\mathbf{d})}$ are concatenated along the channel dimension
and fed into the Latent Diffusion U-Net to predict the noise $\hat{\boldsymbol{\epsilon}}$.
The training objective is the denoising mean squared error:
\begin{equation}
  \mathcal{L}_{\mathrm{Marigold}} =
    \mathbb{E}_{\mathbf{x},\mathbf{d},\boldsymbol{\epsilon},t}
    \bigl\|\boldsymbol{\epsilon} - \hat{\boldsymbol{\epsilon}}_\theta(\mathbf{z}_t^{(\mathbf{d})}, \mathbf{z}^{(\mathbf{x})}, t)\bigr\|^2
\end{equation}

\paragraph{Inference}
At inference time, Marigold adopts DDIM deterministic sampling:
the depth latent vector starts from $\mathbf{z}_T^{(\mathbf{d})} \sim \mathcal{N}(\mathbf{0},\mathbf{I})$,
with the latent code $\mathbf{z}^{(\mathbf{x})}$ of the input image serving as the condition at each denoising step;
after iterative denoising, $\mathbf{z}_0^{(\mathbf{d})}$ is obtained
and decoded back into a depth map $D = \mathcal{D}(\mathbf{z}_0^{(\mathbf{d})})$.
Since depth maps and RGB images share similar local smoothness and spatially correlated structure,
the statistical properties of the VAE latent space apply equally well to the depth modality,
so Marigold can directly reuse the pretrained Stable Diffusion VAE without redesigning the encoder--decoder.

\section{Depth Estimation}
\label{sec:depth}

\hspace{24pt}

As stated in Chapter 1, existing depth estimation models face a significant domain gap on TCP:
whitespace regions are easily misjudged as near-field planes,
and the correspondence between ink-tone variation and depth does not match the priors the models learned from real photographs.
The most direct remedy for this problem would be retraining with a large amount of TCP depth annotations,
but TCP depth datasets are extremely scarce and manual annotation is prohibitively expensive, making this route infeasible in practice.
This study therefore adopts a three-stage human-in-the-loop strategy
to achieve painting-specific depth correction with minimal manual intervention:

\begin{enumerate}
  \item \textbf{Rough depth estimation}:
        a discriminative model (Depth Anything V2) \cite{depthanythingv2} quickly produces a full-image depth preview,
        allowing the user to inspect the depth distribution and locate misjudged regions without spending much time.
        This stage does not pursue precision; its purpose is to provide an initial spatial reference for the manual judgment in the next step.

  \item \textbf{Sparse depth control point annotation}:
        based on the rough depth map, the user only needs to click control points at a few key locations
        (e.g., foreground figures, distant mountains, white space) and assign correct depth values,
        without pixel-wise annotation of the entire image.
        The core assumption of sparse annotation is that, for a domain expert (or a user familiar with the spatial intent of the painting),
        identifying \textit{which regions are misestimated} is far easier than annotating the whole image,
        and these few control points already suffice to guide the model in reconstructing a global depth field consistent with the spatial intent.

  \item \textbf{Depth field refinement with a TCP-aware loss function}:
        the control points provide local numerical depth constraints,
        but optimizing point-wise errors alone cannot guarantee the quality of the overall depth field---
        for example, whitespace regions may still be assigned nonzero depth,
        and unnatural depth discontinuities may appear on both sides of brushstroke boundaries.
        This study proposes \textit{ChinesePaintingLoss},
        which combines a brushstroke-aware smoothness loss with whitespace regularization,
        imposing the above constraints during the denoising optimization of Marigold-DC,
        thereby obtaining a refined depth map $D_{\mathrm{refined}}$ that conforms to the spatial semantics of TCP.
\end{enumerate}

The three steps are logically interlocked:
the rough estimation provides a basis for observation, the sparse annotation injects human spatial intent,
and the TCP-aware loss function ensures the global consistency of the optimization result.
The following sections describe their concrete implementations in turn.

\subsection{Rough Depth Estimation}
\label{subsec:rough}

The rough depth estimation serves only as a reference for the user to inspect the depth distribution and identify problematic regions,
rather than as the final output; hence speed matters more than accuracy at this stage.
This study selects Depth Anything V2 \cite{depthanythingv2} as the rough estimation model.

DA-V2 is a discriminative model that completes inference in a single forward pass,
far faster than diffusion-based methods such as Marigold-DC (which typically require 50 denoising iterations),
and can present a depth preview in real time within the interactive annotation workflow.
Although speed is prioritized, the rich visual priors DA-V2 accumulated from training on large-scale, diverse datasets
already suffice to give the user an informative rough depth distribution.
If the unguided mode of Marigold-DC were used in place of DA-V2,
not only would the iterative cost of diffusion sampling be unsuitable for a quick visual reference,
but the unguided result and the subsequent control-point-guided refined output would merely be different execution paths of the same Marigold-DC,
making it difficult for the user to intuitively judge how much improvement the control point annotation brings.

This study also considered replacing DA-V2 with Depth Anything V3 \cite{depthanything} (hereafter DA-V3).
However, as shown in Table~\ref{tab:da-comparison},
the improvement of DA-V3 over DA-V2 in monocular relative depth manifests mainly on the ETH3D dataset
($\delta_1$ rises from 86.5 to 98.6, about +12 percentage points),
while the differences on KITTI, NYU, SINTEL, and DIODE are rather limited.
ETH3D is a multi-view geometric dataset of real indoor and outdoor scenes;
its scene characteristics differ fundamentally from the stylized two-dimensional spatial expression of TCP,
so this improvement transfers only marginally to the TCP depth estimation task.

\begin{table}[htbp]
  \centering
  \caption{Monocular depth estimation comparison ($\delta_1$ accuracy, \%, higher is better; data source: \cite{depthanything} Table 4)}
  \label{tab:da-comparison}
  \begin{tabular}{lcccccc}
    \hline
    Method & KITTI & NYU & SINTEL & ETH3D & DIODE & Rank \\
    \hline
    Depth Anything V2 & 94.6 & \textbf{97.9} & \textbf{77.2} & 86.5          & 95.2          & 2.60 \\
    VGGT \cite{vggt}  & 91.7 & \textbf{97.9} & 67.9          & 97.5          & 95.3          & 3.75 \\
    Depth Anything V3 & \textbf{95.3} & 97.4 & 75.5 & \textbf{98.6} & \textbf{95.4} & \textbf{2.20} \\
    \hline
    Teacher           & 97.2 & 97.9 & 81.4 & 99.8 & 96.6 & 1.00 \\
    \hline
  \end{tabular}
\end{table}

In the table, Rank denotes the overall ranking of each method across the five datasets, where lower is better;
bold values indicate the best-performing method on each dataset
(Teacher is the teacher model in the DA-V3 paper, serving as an upper-bound reference and not publicly released).

Fundamentally, the core improvement of DA-V3 comes from multi-view geometric understanding---
tasks such as camera pose estimation and 3D reconstruction drive it to learn the geometric structure of real 3D scenes,
whereas TCP suggests spatial layering through artistic devices such as ink-tone variation and the interplay of white space,
rather than being a physical projection of a real 3D scene,
so such multi-view geometric priors offer no direct benefit to TCP depth estimation.
For the scenario addressed in this study,
both DA-V2 and DA-V3 lack corresponding TCP training data
and face a domain gap of comparable magnitude, so the additional generalization capability of DA-V3 provides no substantive advantage here.
Moreover, DA-V3 already assumes the role of metric depth provider in a later stage of the proposed pipeline
(see Section~\ref{sec:pointcloud});
clearly separating the two roles also helps keep the responsibility boundaries between pipeline stages well defined.

This study uses DA-V2 to obtain the rough depth map $D_{\mathrm{rough}}$,
which is presented to the user as the observational basis before sparse control point annotation.

\subsection{Sparse Depth Control Point Annotation}
\label{subsec:dc}

Since depth annotation datasets in the TCP style are extremely scarce,
resolving the domain gap by retraining or fine-tuning existing depth estimation models is infeasible in practice.
This study therefore replaces data-driven training with manual sparse annotation,
directly steering the model's estimates through a small number of user-specified depth control points.
To correct the estimation bias in $D_{\mathrm{rough}}$ caused by the TCP style,
this study develops an interactive annotation tool (Figure~\ref{fig:sparse_depth_hints}).

\begin{figure}[htbp]
  \centering
  \includegraphics[height=5.5cm]{human}
  \hfill
  \includegraphics[height=5.5cm]{sparse_depth_hints}
  \caption{(Left) The input TCP image; (Right) the sparse depth control point annotation interface:
  the user may click an unlimited number of control points anywhere on the image;
  each annotated point is marked with a red box showing its coordinates and initial estimated depth value ($D_{\mathrm{rough}}$),
  and the sliders on the right adjust the depth value of each control point (0 is the farthest, 1 is the nearest).
  Without guidance, the depth estimation of the diffusion model exhibits obvious errors:
  the figure's abdomen should normally protrude more than the head ($D_{\mathrm{rough}}$ closer to 1),
  yet the model output shows the head protruding more than the torso;
  in addition, the background region at the lower left of the image is also assigned a nonzero depth value,
  whereas the background depth should be 0 (the farthest).}
  \label{fig:sparse_depth_hints}
\end{figure}

The user can click control points at locations where the depth judgment is clearly wrong,
and assign correct depth values with sliders (0 is the farthest, 1 is the nearest).
The tool outputs the coordinates and user-specified depth values of the control points,
represented as a sparse depth map $D_{\mathrm{sparse}} \in [0,1]^{H \times W}$
together with a corresponding mask $\mathbf{M} \in \{0,1\}^{H \times W}$:
$\mathbf{M}_{i,j} = 1$ indicates that the pixel carries a user-annotated value,
and $D_{\mathrm{sparse}}$ stores the specified depth of the control point at the corresponding location
(as shown in Figure~\ref{fig:loss_comparison}(b)).

\subsection{Loss Function Design}
\label{subsec:loss}

As described in Chapter 2, supervised models cannot be applied directly to TCP because of the domain gap;
Depth Anything V3, despite its stronger generalization capability, lacks an interface for user intervention,
and the original Marigold cannot accept external constraints at inference time.
Marigold-DC \cite{marigolddc} combines the rich visual priors and zero-shot generalization of diffusion models
with support for guiding depth estimation with sparse depth control points,
making it particularly suitable for the TCP scenario where training data are scarce and local corrections are needed;
this study therefore adopts it as the model for sparse-guided depth estimation.
On top of Marigold's DDIM denoising framework,
Marigold-DC introduces, as its core improvement, a gradient feedback path based on the sparse control points;
the overall guidance procedure is illustrated in Figure~\ref{fig:marigolddc}.
At the end of every denoising step $t$,
Marigold-DC estimates the clean depth latent vector with the Tweedie formula \cite{tweedie}.
The Tweedie formula originates from Bayesian statistics: given a noisy observation,
it directly estimates the posterior mean without performing the full denoising process.
In diffusion models, it takes the form:

\begin{equation}
  \hat{\mathbf{z}}_0^{(\mathbf{d})} =
    \frac{
      \mathbf{z}_t^{(\mathbf{d})} -
      \sqrt{1 - \bar{\alpha}_t}\;
      \hat{\boldsymbol{\epsilon}}_\theta\!\bigl(
        \mathbf{z}_t^{(\mathbf{d})},\, \mathbf{z}^{(\mathbf{x})},\, t
      \bigr)
    }{\sqrt{\bar{\alpha}_t}}
  \label{eq:tweedie}
\end{equation}

where $\bar{\alpha}_t$ is the cumulative noise coefficient of the DDPM forward process,
and $\hat{\boldsymbol{\epsilon}}_\theta$ is the noise predicted by the U-Net at the current time step.
The reason $\hat{\mathbf{z}}_0^{(\mathbf{d})}$ can be computed from $t$ in a single step
lies in the closed-form nature of the DDPM forward noising process:

\begin{equation}
  \mathbf{z}_t = \sqrt{\bar{\alpha}_t}\,\mathbf{z}_0 + \sqrt{1 - \bar{\alpha}_t}\,\boldsymbol{\epsilon},
  \quad \boldsymbol{\epsilon} \sim \mathcal{N}(\mathbf{0}, \mathbf{I})
\end{equation}

that is, the noisy vector at any time step $t$ can be expressed algebraically in terms of the clean vector $\mathbf{z}_0$ and a single noise sample,
with no step-by-step iterative dependency between them.
Therefore, as long as the U-Net gives an accurate prediction
$\hat{\boldsymbol{\epsilon}}_\theta$ of the cumulative noise $\boldsymbol{\epsilon}$ contained in $\mathbf{z}_t$,
$\hat{\mathbf{z}}_0^{(\mathbf{d})}$ can be solved for algebraically,
without stepping back from $t$ through $t{-}1, t{-}2, \ldots, 0$.
Using the U-Net noise prediction, this formula projects the noisy latent vector $\mathbf{z}_t^{(\mathbf{d})}$ back into the clean latent space,
yielding the denoised estimate $\hat{\mathbf{z}}_0^{(\mathbf{d})}$,
which is then converted by the VAE decoder $\mathcal{D}(\cdot)$ into the intermediate depth prediction of the current time step:
$\hat{\mathbf{d}}_t^{(m)} = \mathcal{D}(\hat{\mathbf{z}}_0^{(\mathbf{d})})$.
The loss $\mathcal{L}$ between this prediction and the sparse control points $\mathbf{c}$ (i.e., the depth guidance) is then computed
and backpropagated to $\mathbf{z}_t^{(\mathbf{d})}$,
correcting the latent vector along the direction of the control point constraints before the next denoising step.
By imposing the geometric constraints step by step in this manner, the final depth field
gradually converges to a result consistent with the user annotations without retraining the model.

\begin{figure}[htbp]
  \centering
  \includegraphics[width=0.9\textwidth]{CPL_pipeline}
  \caption{Overall architecture of \textit{ChinesePaintingLoss}.
  The input image $\mathbf{x}$ passes through the Latent Encoder and is denoised by the Marigold-DC backbone;
  at every step the Tweedie formula estimates the intermediate depth $\hat{\mathbf{d}}_t^{(m)}$,
  and the three losses---White Loss (left), Brush Loss (middle), and Metric Loss (right)---are computed simultaneously,
  with backpropagation (red arrows) correcting the latent vector, finally outputting the refined depth map $\hat{\mathbf{d}}^{(m)}$.}
  \label{fig:cpl_pipeline}
\end{figure}

After obtaining the sparse control points, this study uses them as depth guidance and re-runs the Marigold-DC denoising inference with the improved loss function.
The total loss designed in this study, \textit{ChinesePaintingLoss} (shown in Figure~\ref{fig:cpl_pipeline}), comprises one original loss term and two loss terms designed for the characteristics of TCP, defined as:

\begin{equation}
  \mathcal{L}_{\mathrm{final}} =
    \lambda_m \mathcal{L}_{\mathrm{metric}} +
    \lambda_b \mathcal{L}_{\mathrm{brush}} +
    \lambda_w \mathcal{L}_{\mathrm{white}}
  \label{eq:cpl_final}
\end{equation}

All terms are backpropagated jointly at every denoising step.
The three loss terms are described below.

\paragraph{Metric Loss}

The Metric Loss follows the original sparse depth supervision loss of Marigold-DC \cite{marigolddc};
its central purpose is to ensure geometric accuracy---the predicted depth at user-annotated locations must match the control points.
Both L1 and L2 supervision are imposed at the control point locations:

\begin{equation}
  \mathcal{L}_{\mathrm{metric}} =
    \frac{1}{2}\bigl\| D_{\mathrm{pred}}[\mathbf{M}] - D_{\mathrm{sparse}}[\mathbf{M}] \bigr\|_1 +
    \frac{1}{2}\bigl\| D_{\mathrm{pred}}[\mathbf{M}] - D_{\mathrm{sparse}}[\mathbf{M}] \bigr\|_2^2
\end{equation}

where $[\mathbf{M}]$ denotes selecting, via the mask $\mathbf{M}$, the pixel values at all locations with $\mathbf{M}_{i,j}=1$
(i.e., user-annotated positions), which are flattened into a one-dimensional vector before the loss is computed.
The L1 loss is not easily dominated by a few control points with very large deviations, while the L2 loss exerts a stronger corrective force on small deviations;
combined, they effectively suppress deviations of all magnitudes.

As shown in Figure~\ref{fig:loss_comparison}(c),
the Metric Loss alone already makes the depth values at the control point locations converge to the user annotations,
effectively correcting the estimation bias in the rough depth map.
However, the Metric Loss constrains only the annotated locations
and imposes no structural restriction on the regions beyond the control points,
so the heavy ink brushstrokes in TCP are misinterpreted by the model as depth boundaries,
producing unnatural depth ruptures on both sides of the brushstrokes;
to resolve this problem, we introduce the brushstroke-aware smoothness loss.

\paragraph{Brushstroke-aware Smoothness Loss}

The brushstroke-aware smoothness loss proposed in this study (shown in Figure~\ref{fig:bl_pipeline}) is an original design targeting the characteristics of TCP,
built as an improvement upon the traditional image-guided smoothness loss \cite{monodepth}.
The traditional method uses the image gradient magnitude as an exponential decay of the edge weight:
where the gradient is large (brushstroke boundaries) the weight approaches zero, allowing depth discontinuities;
where the gradient is small (uniform regions) the weight approaches one, enforcing strong smoothing.
However, the ink brushstrokes of TCP are heavy, and their luminance gradients are hard to distinguish from those of real object boundaries;
if this mechanism were relied upon entirely, every black brushstroke would be treated as a depth boundary,
forcing the depth map to rupture on both sides of each stroke and producing unnatural geometric segmentation (as shown in Figure~\ref{fig:loss_comparison}(c)).

\begin{figure}[htbp]
  \centering
  \includegraphics[width=\textwidth]{loss_comparison}
  \caption{Comparison of depth estimation results under different loss combinations (from left to right):
  (a) the input TCP image;
  (b) the rough depth map $D_{\mathrm{rough}}$ and the user-annotated sparse control points $D_{\mathrm{sparse}}$;
  (c) the result guided by the Metric Loss only;
  (d) the result of adding the brushstroke-aware smoothness loss $\mathcal{L}_{\mathrm{brush}}$ on top of (c);
  (e) the result of further adding the whitespace regularization $\mathcal{L}_{\mathrm{white}}$ on top of (d),
  i.e., the complete \textit{ChinesePaintingLoss}.
  From (c) to (d), the depth ruptures at the brushstrokes are alleviated and the overall depth becomes smoother;
  from (d) to (e), the depth of the whitespace regions (red regions) is correctly pushed down to the far range.}
  \label{fig:loss_comparison}
\end{figure}

To this end, this study sets a minimum connectivity lower bound $w_{\min}$ on the edge weight,
ensuring that even at brushstroke boundaries with extremely large gradients, the smoothing pressure never drops to zero;
the image gradient $\nabla I^{\mathrm{gray}}$ is computed with the Sobel filter:

\begin{equation}
  w_{i,j} = \max\!\Bigl( e^{-\alpha \|\nabla I^{\mathrm{gray}}\|_{i,j}},\; w_{\min} \Bigr)
\end{equation}

where $e^{-\alpha \|\nabla I^{\mathrm{gray}}\|}$ is the original edge weight,
which decays to nearly 0 at brushstroke boundaries (where $\|\nabla I^{\mathrm{gray}}\|$ is large),
while $w_{\min} = 0.3$ serves as the lower bound, guaranteeing that 30\% of the smoothing pressure is retained there,
preventing the depth map from forming non-physical hard ruptures along brushstrokes.
The complete brushstroke smoothness loss proposed in this study is:

\begin{equation}
  \mathcal{L}_{\mathrm{brush}} =
    \frac{1}{HW} \sum_{i,j} w_{i,j}
    \Bigl(
      \bigl|\partial_x D_{\mathrm{pred}}\bigr|_{i,j} +
      \bigl|\partial_y D_{\mathrm{pred}}\bigr|_{i,j}
    \Bigr)
\end{equation}

where $\alpha = 1.0$.
$I^{\mathrm{gray}} \in [0,1]$ is the grayscale-normalized image, converted with the standard luminance weighting formula:

\begin{equation}
  I^{\mathrm{gray}} = 0.299\,R + 0.587\,G + 0.114\,B
\end{equation}

\begin{figure}[htbp]
  \centering
  \includegraphics[width=0.9\textwidth]{BL_pipeline}
  \caption{Computation procedure of the brushstroke-aware smoothness loss $\mathcal{L}_{\mathrm{brush}}$. The spatial gradient $|\nabla\hat{\mathbf{d}}_t^{(m)}|_1$ of the predicted depth $\hat{\mathbf{d}}_t^{(m)}$ is taken; the input grayscale image $I^{\mathrm{gray}}$ passes through the Sobel filter to obtain the image gradient $\nabla I^{\mathrm{gray}}$, which is then converted into the edge weights $w_{i,j}$; the product of the two gives the per-pixel loss contribution, automatically suppressing the penalty along brushstroke contours (low weight) while maintaining depth smoothness in flat regions (high weight).}
  \label{fig:bl_pipeline}
\end{figure}

As shown in Figure~\ref{fig:loss_comparison}(d), after adding $\mathcal{L}_{\mathrm{brush}}$,
the depth ruptures at the brushstrokes are clearly improved and the overall depth field becomes continuous.
However, another type of TCP-specific estimation bias remains in the depth map:
large whitespace regions are assigned depth values close to the foreground,
whereas in the visual language of TCP, white space represents the distant sky or the void, and its depth should approach the far end.
Since the proposed $\mathcal{L}_{\mathrm{brush}}$ only constrains spatial smoothness in high-gradient regions
and imposes no semantic restriction on the absolute depth values of whitespace regions,
this study designs a separate regularization term for the whitespace semantics.

\paragraph{Whitespace Regularization}

TCP expresses distant scenery, sky, or the void with large areas of white space,
a visual language by which traditional Chinese painting interprets \textit{infinite distance} through \textit{nothingness},
or simply the background beyond the main subject.
However, depth estimation models are trained on real photographs,
in which large bright regions usually correspond to nearby planes such as white walls or tabletops,
so the models tend to misjudge TCP white space as foreground, assigning depth values close to 1 (nearest).
Such misjudgments cannot be corrected by sparse control points alone---
the user cannot possibly annotate a control point on every patch of white space,
so a semantic constraint must be added to the loss function,
letting the model automatically push high-luminance regions toward the far field during optimization.

The proposed approach proceeds in two steps (shown in Figure~\ref{fig:wl_pipeline}):

\textbf{Step 1: whitespace detection.}
The Sigmoid function is defined as:
\begin{equation}
  \sigma(x) = \frac{1}{1 + e^{-x}}
\end{equation}
Its output range is $(0, 1)$, mapping any real number to a continuous probability value.
It is applied to the grayscale value $I^{\mathrm{gray}}_{i,j}$ of each pixel
to compute a continuous whitespace probability $M^{\mathrm{white}}_{i,j}$:

\begin{equation}
  M^{\mathrm{white}}_{i,j} = \sigma\!\bigl((I^{\mathrm{gray}}_{i,j} - \tau) \cdot 20\bigr)
\end{equation}

where $\tau = 0.85$ is the whitespace luminance threshold, and the constant $20$ controls the steepness of the transition.
When the grayscale value is clearly above $\tau$ (close to pure white), $M^{\mathrm{white}}_{i,j} \to 1$ and the pixel is judged as white space;
when the grayscale value is below $\tau$ (ink present), $M^{\mathrm{white}}_{i,j} \to 0$ and no constraint is imposed.
The sigmoid is used instead of a hard threshold
(i.e., the indicator function $\mathbf{1}[I^{\mathrm{gray}} > \tau]$, which outputs 1 when the luminance exceeds $\tau$ and 0 otherwise)
because a hard threshold makes a binary decision:
pixels falling on either side of the threshold receive either the full constraint or none at all, no matter how tiny the difference.
However, TCP abounds in ambiguous \textit{semi-whitespace} regions such as light ink wash blending and uneven rice-paper background tone;
a binary cut would jag the mask edges and make the result extremely sensitive to the choice of threshold.
The sigmoid instead outputs a weight that increases smoothly with luminance:
pixels close to pure white receive nearly the full constraint,
while light-ink regions receive only a proportionally mild constraint,
so that the strength of the penalty matches the degree to which a pixel \textit{looks like white space},
making it more robust to ambiguous regions and to the choice of threshold.

\textbf{Step 2: penalizing depth deviation in whitespace regions.}
Using $M^{\mathrm{white}}_{i,j}$ as the weight,
the deviation of the predicted depth from the farthest end $d_{\mathrm{far}} = 0$ is penalized:

\begin{equation}
  \mathcal{L}_{\mathrm{white}} =
    \frac{1}{HW} \sum_{i,j} M^{\mathrm{white}}_{i,j} \cdot D_{\mathrm{pred},i,j}^{\,2}
\end{equation}

If the model correctly predicts white space as far field ($D_{\mathrm{pred}} \approx 0$), the loss contribution is close to zero;
if it is misjudged as foreground ($D_{\mathrm{pred}} \approx 1$), the loss is amplified quadratically, imposing a strong penalty.

\begin{figure}[htbp]
  \centering
  \includegraphics[width=0.85\textwidth]{WL_pipeline}
  \caption{Computation procedure of the whitespace regularization $\mathcal{L}_{\mathrm{white}}$. The grayscale image $I^{\mathrm{gray}}$ passes through the Sigmoid function with threshold $\tau = 0.85$ to obtain the whitespace mask $M^{\mathrm{white}}$, which is then multiplied pixel-wise with the squared predicted depth $\left(\hat{\mathbf{d}}_t^{(m)}\right)^{\!2}$, penalizing nonzero depth in whitespace regions and pushing the model to move high-luminance regions toward the farthest end ($D \to 0$).}
  \label{fig:wl_pipeline}
\end{figure}

As shown in Figure~\ref{fig:loss_comparison}(e), after adding this term, the depth of the whitespace regions (inside the red boxes) is clearly pushed down.

\paragraph{Weight Settings and Output}
The weight coefficients of the loss terms (Eq.~(\ref{eq:cpl_final})) are set to:
$\lambda_m = \dfrac{2}{3}$, $\lambda_b = \dfrac{1}{6}$, $\lambda_w = \dfrac{1}{6}$.

After all losses converge, the final depth prediction output by the denoising process is denoted as the refined depth map
$D_{\mathrm{refined}} \in [0,1]^{H \times W}$, which is used for the subsequent unprojection.

\section{Point Cloud Generation and Rendering}
\label{sec:pointcloud}

The point cloud reconstruction in this section takes the refined depth map $D_{\mathrm{refined}}$ of a single TCP image as its sole input,
converting it into a three-dimensional point cloud via single-view unprojection,
without involving multi-view geometric reconstruction or 3D model construction.

Intuitively, given a depth map $D$ and a camera intrinsic matrix $\mathbf{K}$:

\begin{equation}
  \mathbf{K} =
  \begin{pmatrix} f_x & 0 & c_x \\ 0 & f_y & c_y \\ 0 & 0 & 1 \end{pmatrix}
\end{equation}

where $f_x, f_y$ are the focal lengths and $(c_x, c_y)$ is the principal point,
every pixel $(u,v)$ can be unprojected directly into three-dimensional space:

\begin{equation}
  \begin{pmatrix} X \\ Y \\ Z \end{pmatrix} =
  D(u,v) \cdot \mathbf{K}^{-1}
  \begin{pmatrix} u \\ v \\ 1 \end{pmatrix}
\end{equation}

However, $D_{\mathrm{refined}}$ suffers from two fundamental limitations that prevent it from serving directly as the input to this formula:

\begin{enumerate}
  \item \textbf{Lack of metric scale}: $D_{\mathrm{refined}}$ is an affine-invariant relative depth
        with values in $[0,1]$, not true depth in meters.
        A point cloud unprojected from relative depth has the correct geometric shape but an unknown absolute scale;
        it cannot be matched to a real camera trajectory, and the parallax appears distorted during rendering.
  \item \textbf{Lack of camera intrinsics}: $\mathbf{K}^{-1}$ converts pixel coordinates $(u,v)$ back into direction vectors in the camera coordinate system,
        which is a necessary condition for computing 3D coordinates.
        Marigold-DC performs inference via diffusion denoising; it does not estimate camera geometric parameters and cannot output $\mathbf{K}$.
\end{enumerate}

This study therefore introduces Depth Anything V3 \cite{depthanything}
to resolve the two limitations above and obtain an accurate three-dimensional point cloud.
DA-V3 is trained with supervision on large-scale, diverse data; it can estimate metric depth directly from the input image
and internally estimates the camera intrinsic matrix $\mathbf{K}$;
its DA3\_ToPointCloud module combines the two computations and outputs a 3D point cloud directly,
so the user does not need to manually extract or pass $\mathbf{K}$ during point cloud generation.
This study feeds the grayscale visualization of $D_{\mathrm{refined}}$ into DA-V3,
letting it output a metric-scaled three-dimensional point cloud
while preserving the spatial structure refined by Marigold-DC, for subsequent rendering by Uni3C.

The output point cloud is colored with the pixel colors of the original input TCP image,
simultaneously recording the geometric position and color information of every point in 3D space;
it can be rendered along a predefined camera trajectory using standard computer-graphics rasterization,
producing a preliminary 2.5D point-cloud video with true parallax.
The rendered result contains occlusion holes exposed by viewpoint shifts,
but it provides explicit spatial-geometric conditions for the subsequent video diffusion model.

\section{Video Diffusion Generation}
\label{sec:video}

\begin{figure}[htbp]
  \centering
  \includegraphics[width=\textwidth]{TCP_uni3c}
  \caption{Overview of the complete TCP-Uni3C inference pipeline.
  The upper half is \textit{ChinesePaintingLoss} (White Loss, Brush Loss, Metric Loss),
  which guides Marigold-DC to output the refined depth map $D_{\mathrm{refined}}$,
  from which Depth Anything V3 unprojects a geometrically accurate point cloud.
  The lower half is the video generation stage:
  this study designs a circular camera trajectory and renders the geometrically accurate point cloud into a frame-by-frame point-cloud video $V_{\mathrm{pcd}}$ with hole masks (Masks),
  which, together with positive and negative text prompts designed for the static-scene characteristics of TCP,
  serve as the guidance conditions of Uni3C \cite{uni3c},
  driving Wan-DiT \cite{wan2025} to generate a 2.5D video with parallax and consistent style.}
  \label{fig:tcp_uni3c}
\end{figure}

Although the TCP point cloud rendering improved in this study provides correct parallax, two fundamental limitations remain.
First, when the camera translates or rotates, background regions originally occluded by the foreground are gradually exposed,
and since these regions have no corresponding geometric data in the point cloud, holes appear in the rendered result.
Second, rasterization can only project discrete points onto the image plane;
it cannot reproduce the texture details of ink brushwork or the spatial atmosphere of white space in TCP,
so the directly rendered frames lack visual quality.
If the holes were filled with traditional interpolation or image inpainting,
the results, lacking semantic scene understanding, would often clash with the TCP style;
especially when the 2.5D viewpoint shift is large, the hole regions expand accordingly,
and the drawbacks of traditional methods become more pronounced, easily producing blurry gradations or discontinuous visual artifacts, as shown in Figure~\ref{fig:traditional_comp}.


Video diffusion models, by contrast, inherently possess two complementary capabilities:
first, they fill occlusion holes with learned visual priors,
generating content that is semantically and stylistically consistent with the surrounding frames;
second, they maintain visual consistency along the temporal dimension of the entire video,
avoiding the flickering caused by frame-by-frame independent inpainting.
However, pure video diffusion models trained mainly on real-scene footage lack 3D geometric constraints;
the generated camera motion can hardly correspond precisely to the predefined trajectory,
easily producing geometrically inconsistent hallucinations during viewpoint changes.

This study therefore introduces Uni3C \cite{uni3c},
which uses the point cloud as a geometric guidance condition to drive a video diffusion model in generating 2.5D videos,
and on this basis proposes an improved architecture for the TCP domain---
Traditional Chinese Painting Uni3C (TCP-Uni3C).
TCP-Uni3C uses the geometrically accurate point-cloud video $V_{\mathrm{pcd}}$ obtained by unprojecting $D_{\mathrm{refined}}$,
together with text prompts designed for static TCP scenes and a circular camera trajectory,
to jointly drive the denoising process of Wan-DiT \cite{wan2025}, as shown in Figure~\ref{fig:tcp_uni3c}.

\begin{figure}[htbp]
  \centering
  \includegraphics[width=\textwidth]{tradition_comp}
  \caption{Comparison of traditional inpainting methods under different viewpoint offsets.
    The thumbnail on the left is the original TCP reference frame; the grid on the right shows four viewpoint offsets: Frame 3, 6, 9, and 12.
    The top row shows the point-cloud rendered frames (gray regions are holes caused by the 2.5D offset),
    the middle row shows the 2.5D video generation results produced by this study,
    and the bottom row shows the PatchMatch block-matching inpainting results.
    It can be seen that the traditional method produces obvious blur and artifacts in the hole regions
    and cannot preserve the sharp line quality of TCP brushwork or the atmosphere of its white space.}
  \label{fig:traditional_comp}
\end{figure}

If the original Uni3C were used directly, its built-in Depth-Pro \cite{depthpro} would be responsible for monocular depth estimation of the input image,
unprojecting the estimate into a point cloud and rasterizing it along the camera trajectory into the frame-by-frame video $V_{\mathrm{pcd}}$.
However, as described in Section~\ref{sec:depth}, Depth-Pro is trained on real-scene photographs;
faced with the whitespace composition, ink textures, and abstract spatial relations unique to TCP, it can hardly produce geometrically reliable depth estimates, as shown in Figure~\ref{fig:uni3c_comp},
and the resulting point cloud is not accurate enough to drive a 2.5D animation with realistic parallax.

\begin{figure}[htbp]
  \centering
  \includegraphics[width=\textwidth]{uni3c_comp}
  \caption{Comparison of the generation results of the original Uni3C and the proposed TCP-Uni3C.
    The leftmost image (first frame) is the input TCP image, which also serves as the first frame of the video;
    the middle columns show the point-cloud video $V_{\mathrm{pcd}}$ and the final generated video of the original Uni3C,
    and the right columns show the corresponding $V_{\mathrm{pcd}}$ and generated video of TCP-Uni3C.
    It can be seen that the original Uni3C, driven by the flattened depth estimated by Depth-Pro,
    produces the effect of a camera orbiting around a picture frame, lacking true depth parallax;
    in contrast, TCP-Uni3C, driven by the refined depth $D_{\mathrm{refined}}$,
    generates a 2.5D video with genuine parallax.}
  \label{fig:uni3c_comp}
\end{figure}

\paragraph{Camera Intrinsic Matrix}

TCP-Uni3C therefore directly replaces the output of Depth-Pro with the accurate point cloud generated by the aforementioned pipeline,
retaining only the focal length $f_{\mathrm{px}}$ that Depth-Pro estimates for the reference image,
which is used to construct the camera intrinsic matrix required for rendering
\begin{equation}
\mathbf{K} = \begin{pmatrix} f_{\mathrm{px}} & 0 & w/2 \\ 0 & f_{\mathrm{px}} & h/2 \\ 0 & 0 & 1 \end{pmatrix}
\end{equation}
where pixels are assumed to be square, so the horizontal and vertical focal lengths are equal ($f_{\mathrm{px}} = f_{\mathrm{py}}$) and both denoted by $f_{\mathrm{px}}$;
$(w/2,\, h/2)$ is the principal point of the image,
and $w$, $h$ are the image width and height (in pixels), respectively.
The initial distance $r$ between the camera and the point cloud is derived from the principle of similar triangles (as shown in Figure~\ref{fig:camera_geometry}),
computed from the 3D height $H_{\mathrm{pcd}}$ of the point cloud, the focal length $f_{\mathrm{px}}$, and the image height $h$:
\begin{equation}
r = \frac{H_{\mathrm{pcd}} \cdot f_{\mathrm{px}}}{h}
\label{eq:camera_dist}
\end{equation}
ensuring that the projected point cloud exactly fills the first frame.
However, this distance sometimes places the camera too far from the point cloud, so that when the camera moves, overly large hole regions appear in the background,
degrading the filling quality of the subsequent diffusion model, as shown in Figure~\ref{fig:camera_distance}.
This study therefore manually shortens the camera distance when necessary (i.e., adjusts the scaling coefficient of $r$),
bringing the camera closer to the point cloud as a whole to reduce the hole area and improve the quality of the final generated video.

\begin{figure}[htbp]
  \centering
  \includegraphics[width=0.82\textwidth]{camera_geometry}
  \caption{Geometric derivation of the camera distance $r$.
  The camera, the image plane (height $h$, at distance $f_{\mathrm{px}}$ from the camera), and the point cloud (3D height $H_{\mathrm{pcd}}$, at distance $r$ from the camera)
  form similar triangles, and the initial camera distance follows directly from the proportion $H_{\mathrm{pcd}}/r = h/f_{\mathrm{px}}$.}
  \label{fig:camera_geometry}
\end{figure}

\begin{figure}[htbp]
  \centering
  \includegraphics[width=0.9\textwidth]{camera_distance}
  \caption{Effect of the camera distance on the hole regions.
  The top row shows the results of rendering the point cloud along the configured camera trajectory,
  and the bottom row shows the results after filling by the diffusion model.
  (a) When the camera is too far from the point cloud, large background holes are exposed after the viewpoint shift;
  (b) after shortening the camera distance, the hole area shrinks noticeably, facilitating the filling by the diffusion model.}
  \label{fig:camera_distance}
\end{figure}

\paragraph{Camera Extrinsics and Trajectory}

Once the camera distance $r$ is determined, the camera rotation and position of every frame---i.e., the extrinsic matrix $[\mathbf{R}_i \mid \mathbf{t}_i]$---must be further defined.
This study adopts a circular trajectory: taking the initial camera view $\mathbf{C}_0$ as the reference,
each frame is given an azimuth offset $\Delta\phi_i$ (horizontal direction)
and an elevation offset $\Delta\theta_i$ (vertical direction).
The video has $F$ frames in total, of which the first frame applies no angular offset and keeps the initial camera view $\mathbf{C}_0$ (corresponding to the original composition of the input painting);
the angular offsets of the remaining $N = F-1$ frames (frames $2$ through $F$ of the video) are determined by the parameter $t_i = \frac{2\pi i}{N}$ ($i = 0, \ldots, N-1$):

\begin{equation}
  \Delta\phi_i   =  A_\phi \sin(t_i), \qquad
  \Delta\theta_i = -A_\theta \cos(t_i)
\end{equation}

where $A_\phi$ and $A_\theta$ are the angular amplitudes in the horizontal and vertical directions, respectively,
which can be set independently according to the composition of the painting and the desired parallax strength;
this study by default sets the amplitudes within the range of $3^\circ$ to $15^\circ$, adjusted per painting,
i.e., the maximum camera offset relative to the initial position lies between $\pm 3^\circ$ and $\pm 15^\circ$,
and the overall motion range covers $6^\circ$ to $30^\circ$.
The amplitude cap is deliberately kept within $15^\circ$ because, when the camera offset angle is too large,
the exposed hole regions in the point-cloud-rendered $V_{\mathrm{pcd}}$ expand accordingly,
increasing the proportion of content the diffusion model must complete on its own and degrading the quality of the final generated video;
within this range, the hole area remains within what the diffusion model can complete stably.
There is a $90^\circ$ phase difference between the horizontal and vertical offsets ($\sin$ versus $\cos$),
so that together they trace a closed circular path in the $(\Delta\phi,\,\Delta\theta)$ space
(an ellipse if the two amplitudes are set to different values),
ensuring that the camera produces smooth parallax motion in both the horizontal and vertical directions simultaneously;
it is worth noting that at $t_0 = 0$ the offset is $(0,\,-A_\theta)$,
i.e., the trajectory frames do not start from $\mathbf{C}_0$ but enter the circular path at a point $A_\theta$ above the initial view;
furthermore, since the parameter samples are uniformly distributed over $[0,\,2\pi)$ excluding the endpoint,
the last frame lands exactly one step before the trajectory entrance, so the entire orbiting motion joins smoothly end to end.
The 3D visualization of the complete trajectory is shown in Figure~\ref{fig:camera_traj}.

\begin{figure}[htbp]
  \centering
  \includegraphics[width=0.85\textwidth]{camera_traj}
  \caption{3D visualization of the circular camera trajectory.
    The green dot (start) is the initial camera position (corresponding to $\mathbf{C}_0$), located at the center of the circular trajectory and corresponding to the first frame of the video;
    the blue ring is the closed circular trajectory formed by the angular offsets $(\Delta\phi,\,\Delta\theta)$,
    and the gray cones are the camera frustums sampled along the trajectory;
    the red square (end) is the camera position of the last frame, located near the trajectory entrance.
    On the right is the three-dimensional point cloud unprojected from the TCP image.}
  \label{fig:camera_traj}
\end{figure}

The camera-to-world matrix of each frame is obtained by composing the corresponding rotation matrices:

\begin{equation}
  \mathbf{C}_i = \mathbf{R}_\phi(\Delta\phi_i)\,\mathbf{R}_\theta(\Delta\theta_i)\,\mathbf{C}_0
\end{equation}

where $\mathbf{R}_\phi$ is the rotation matrix about the $y$-axis and $\mathbf{R}_\theta$ is the rotation matrix about the $x$-axis;
the extrinsic matrix is taken from the first three rows of its inverse: $[\mathbf{R}_i \mid \mathbf{t}_i] = \bigl[\mathbf{C}_i^{-1}\bigr]_{3\times4}$, where the subscript $3\times4$ denotes truncation to the first three rows.
The extrinsic matrix of each frame, together with the intrinsic matrix $\mathbf{K}$, determines the rasterized projection of the point cloud for that frame,
constituting the subsequent geometric guidance conditions.

The point-cloud guidance mechanism of Uni3C simultaneously solves the hole-filling problem:
for pixels covered by the point-cloud projection, the model generates under the geometric constraints of the point cloud, ensuring consistency with the 3D structure;
for the hole regions exposed by the viewpoint shift, where the point cloud lacks corresponding geometric data,
Wan-DiT instead fills them automatically using the visual appearance prior of the first frame and the semantic guidance of the text prompts,
generating content that is stylistically and semantically consistent with the surrounding frames,
avoiding the unnatural results that traditional image inpainting produces in the absence of semantic understanding.

\paragraph{Text Prompts for 2.5D Visual Presentation}

Without modifying the main architecture of Uni3C,
this study designs dedicated text prompts for the underlying Wan2.1 I2V model, tailored to the characteristics of TCP animation generation.
The positive prompt is set to:
\begin{quote}
\textit{``This video describes a slow and stable camera movement with a static 3D object.''}
\end{quote}
The negative prompt is set to:
\begin{quote}
\textit{``object movement''}
\end{quote}
The positive prompt guides the model to interpret the scene as camera motion rather than object motion,
ensuring that the content generated when filling the holes matches the semantic expectation of a static scene under viewpoint changes;
the negative prompt explicitly suppresses the generation of object-motion signals,
preventing the model from introducing object-displacement hallucinations inconsistent with TCP semantics when filling the occluded regions,
thereby preserving the static nature of the figures and scenery in the painting.
Finally, Uni3C outputs a smooth, complete 2.5D video with realistic parallax.
